% !TeX spellcheck = it_IT

\documentclass[handout]{beamer}

\usepackage[utf8]{inputenc}
\usepackage[italian]{babel}
\usepackage{multicol}
\usetheme{metropolis}

\usepackage{listings}

\usepackage{tikz}
\usetikzlibrary{shapes,arrows,positioning}

\tikzstyle{startstop} = [ellipse, draw, fill=bg, text=structure, text width=3em, text badly centered, inner sep=0pt]
\tikzstyle{block} = [rectangle, draw, fill=bg, text=structure, text width=5em, text centered]
\tikzstyle{decision} = [diamond, aspect=2, draw, fill=bg, text=structure, text width=3em, text badly centered, inner sep=0pt]
\tikzstyle{input} = [rectangle, draw, fill=bg, text=structure, text width=4em, text centered, rounded rectangle, rounded rectangle east arc=none]
\tikzstyle{output} = [rectangle, draw, fill=bg, text=structure, text width=4em, text centered, rounded rectangle, rounded rectangle west arc=none]

\tikzstyle{line} = [draw, -latex']


\title{Informatica B - Esercitazione 8}
\subtitle{Matlab - Struct e funzioni}

\author{Stefano Cereda\\
		stefano1.cereda@mail.polimi.it
	}
\date{28/11/2017}
\institute[PoliMi]{\vspace{0.5cm}\centering Politecnico di Milano \\ \vspace{0.2cm}
	\includegraphics[width=0.2\linewidth]{./logopolimi}}
%\logo{\includegraphics[width=15mm]{./logopolimi}}

\setbeamercovered{invisible}

\makeindex

\begin{document}
\begin{frame}
	\maketitle
\end{frame}


\begin{frame}{Esercizio 2 - caffè - non fatto}
In un dipartimento ci sono 9 persone, ognuna nell'ultima settimana ha
consumato un certo numero di caffè memorizzati in una matrice "caffe" data
(dove le righe rappresentano le persone, mentre le colonne i giorni).
\begin{enumerate}
\item Si vuole sapere quanto deve pagare ogni persona sapendo che ogni caffè
ha un costo di 0,25 euro.
\item Si supponga che vi sia un costo fisso di noleggio della macchinetta per
fare il caffè pari a 10 euro a settimana.
\begin{enumerate}
\item  La metà di tale costo è divisa equamente tra i bevitori (coloro che
hanno consumato almeno un caffè);
\item l'altra metà tra i bevitori che ne hanno consumato di più rispetto
alla media.
\end{enumerate}
\end{enumerate}

Si dica quanto deve pagare ogni persona.
\end{frame}

\begin{frame}{Esercizio 3 - Paghe - non fatto}
	Data la seguente tabella:
	
\begin{tabular}{|c|c|c|c|c|c|}
	\hline
	& op1 & op2 & op3 & op4 & op5 \\
	\hline
	Paga oraria & 5 & 5.5 & 6.5 & 5 & 6.25 \\ 
	\hline
	Ore settimanali & 40 & 43 & 37 & 50 & 45 \\
	\hline
	Pezzi prodotti & 1000 & 1100 & 1000 & 1200 & 1100 \\
	\hline
\end{tabular} 

Scrivere uno script Matlab che risponda alle seguenti domande:
\begin{enumerate}
\item Quanto guadagna ogni operaio?
\item Qual è il salario totale di tutti gli impiegati?
\item Quanti pezzi vengono prodotti?
\item Qual è il costo medio di un pezzo?
\item Quante ore occorrono in media per un pezzo?
\item Qual è l’operaio più efficiente?
\end{enumerate}
\end{frame}


\begin{frame}[allowframebreaks]{Esercizio 4 - Archivio mammiferi - dati strutturati}

	Si vogliono rappresentare informazioni relative ai mammiferi ospitati in un
	parco protetto.
	
	In particolare, per ogni mammifero si rappresentano i seguenti dati:\\
	NOME TIPO, TIPO ALIMENTAZIONE, PESO MEDIO, insieme di ESEMPLARI
	
	Dove ogni esemplare è caratterizzato da:\\
	CODICE ESEMPLARE, ANNO DI NASCITA, NUMERO CUCCIOLI
	
	\begin{enumerate}
	\item Si crei una variabile mammifero che contenga le informazioni del
	seguente mammifero:\\
	NOME TIPO: Elefante\\
	TIPO ALIMENTAZIONE: Erbivoro\\
	PESO MEDIO: 2500\\
	ESEMPLARE 1: (Codice: 3, Anno: 1985, Cuccioli: 3)\\
	ESEMPLARE 2: (Codice: 5, Anno: 1989, Cuccioli: 0)\\
	ESEMPLARE 3: (Codice: 8, Anno: 1982, Cuccioli: 0)\\
	
	\item Si scriva il codice che costruisce, a partire dalla variabile creata al punto
	precedente, l'archivio degli esemplari di elefante sterili sapendo che un
	elefante si dice sterile se è nato prima del 1990 e non ha concepito
	nessun cucciolo.
\end{enumerate}

	Si stampi su schermo l'elenco dei codici degli elefanti sterili.
	
\end{frame}

\begin{frame}{Esercizio 5 - funzione split}
Scrivere una funzione (split) che data una variabile di ingresso $m$ (scalare, vettore o
matrice di numeri) ed un numero $x$ restituisce due vettori $vMin$ e $vMax$. $vMin$
contiene tutti gli elementi di $m$ minori di $x$, mentre $vMax$ gli elementi maggiori
di $x$.
\end{frame}

\begin{frame}[allowframebreaks]{Esercizio 6 - codice IBAN}
Un codice IBAN è una sequenza di 27 caratteri alfanumerici così composta:
\begin{itemize}
\item 2 caratteri maiuscoli rappresentanti la sigla internazionale della Nazione
\item 2 cifre (CIN Europeo)
\item 1 carattere maiuscolo (CIN Italiano)
\item 5 cifre (ABI)
\item 5 cifre (CAB)
\item 12 cifre (Numero di Conto)
\end{itemize}
\begin{tabular}{|c|c|c|c|c|c|}
	\hline
	naz & cin\_e & cin\_I & abi & cab	&	conto \\
	\hline
	IT&02&L&12345&12345&123456789012\\
	\hline
\end{tabular}


Si scrivano le seguenti funzioni:
\begin{itemize}
\item \emph{removeSpaces}, che prenda in ingresso una stringa strIn e restituisca una stringa strOut contenente tutti i caratteri di strIn ad esclusione degli spazi.
\item \emph{isAlpha}, che prenda in ingresso una stringa str e restituisca 1 quando str contiene esclusivamente lettere maiuscole, o 0 altrimenti.
\item \emph{isDigit}, che prenda in ingresso una stringa str e restituisca 1 quando str contiene esclusivamente caratteri corrispondenti a cifre, o 0 altrimenti.
\end{itemize}

Si utilizzino tali funzioni per scrivere un’ulteriore funzione \emph{checkIBAN} che riceve come parametro una stringa e restituisce 1 se la stringa inserita, una volta rimossi gli spazi, corrisponde ad un codice IBAN che rispetta lo schema sopra illustrato, o 0 altrimenti.
\end{frame}

	
\begin{frame}[allowframebreaks]{Esercizio 7 - biblioteca}
Si vuole automatizzare la gestione di una biblioteca di un piccolo paese.

Come primo passo si vuole gestire una matrice di dimensioni $N \times 2$ contenente nella
prima colonna un identificativo numerico e nella seconda colonna un numero di
tessera.

Immaginate l’identificativo numerico come un codice interno alla biblioteca mentre il
numero di tessera come il numero di una tessera esistente appartenente (ad esempio
il vostro codice persona del politecnico).

\begin{tabular}{|c|c|}
	\hline
	ID	&	Codice Persona \\
	\hline
	1	&	123456789	\\
	\hline
	2	&	987654321	\\
	\hline
	...	&	... \\
	\hline
\end{tabular} 


Si forniscano le funzioni necessarie all'inserimento di una nuova tessera e alla ricerca
dell'identificativo numerico dato un numero di tessera (l’identificativo numerico è un
numero crescente che parte da 1).

Funzione \emph{inserisci}:\\
\emph{input}: matrice Soci e nuovo numero di tessera\\
\emph{output}: la matrice Soci

Funzione \emph{cerca}:\\
\emph{input}: matrice Soci e numero di tessera da cercare\\
\emph{output}: l’identificativo numerico associato al numero tessera
\end{frame}


\end{document}